% Figure: normal-incidence absorption coefficient versus standing-wave
% ratio in an impedance tube, alpha = 1 - ((s-1)/(s+1))^2.
\begin{figure}[htbp]
\centering
\begin{tikzpicture}
\begin{axis}[
  width=12cm, height=7.5cm,
  xlabel={standing-wave ratio $s = p_{\max}/p_{\min}$},
  ylabel={absorption coefficient $\alpha$},
  xmode=log,
  xmin=1, xmax=100, ymin=0, ymax=1.05,
  axis lines=left,
  domain=1:100, samples=300,
  every axis x label/.style={at={(current axis.right of origin)}, anchor=west},
  every axis y label/.style={at={(current axis.above origin)}, anchor=south},
  grid=major, grid style={enggray!20},
  tick label style={font=\scriptsize},
  label style={font=\small},
]
% alpha = 1 - |R|^2, with |R| = (s-1)/(s+1).
\addplot[engblue, very thick] {1 - ((x-1)/(x+1))^2};
% mark the perfect-absorber limit s -> 1
\addplot[enggray, thin, dashed] coordinates {(1,0) (1,1.05)};
\node[enggray, font=\scriptsize, anchor=south west] at (axis cs:1.05,0.9)
  {$s\to1$: perfect absorption};
% mark a worked point s = 3 -> alpha = 0.75
\addplot[engred, only marks, mark=*, mark size=1.5pt] coordinates {(3,0.75)};
\node[engred, font=\scriptsize, anchor=west] at (axis cs:3.3,0.72)
  {$s=3,\ \alpha=0.75$};
\end{axis}
\end{tikzpicture}
\caption[Absorption coefficient from the standing-wave ratio]{Normal-
incidence absorption coefficient as a function of the standing-wave ratio
$s=p_{\max}/p_{\min}$ measured in an impedance (Kundt) tube. A pure tone
driven down the tube reflects from the test sample at the far end, and
the incident and reflected waves form a standing-wave pattern whose
maximum-to-minimum pressure ratio is $s$. The reflection-coefficient
magnitude is $|R|=(s-1)/(s+1)$, so the absorption coefficient is
$\alpha=1-|R|^2$. A perfectly absorbing sample gives no reflected wave,
a flat field, and $s=1$ with $\alpha=1$; a rigid reflector gives
$s\to\infty$ and $\alpha=0$. The marked point $s=3$ returns
$\alpha=0.75$. The tube method measures the normal-incidence coefficient
on a small sample, which differs from the random-incidence Sabine
coefficient of the reverberation-room method; the two are complementary
standardised tests, the tube fast and material-economical, the chamber
representative of a diffuse field.}
\label{fig:vol03ch07:impedance-tube}
\end{figure}
